\chapter{Derivation R: The Resurgence Cancellation (Non-Perturbative Ambiguity)}
\label{ch:derivation_r_resurgence}

\section{Context: The Un-Summability of Perturbation Theory}
A rigorous definition of Yang-Mills theory must address the fact that the perturbative expansion in the coupling constant $g^2$ is divergent. The coefficients $a_n$ of the series $\sum a_n (g^2)^n$ grow factorially as $n!$, rendering the series non-Borel summable.
This leads to an imaginary ambiguity in the summation, which must be exactly cancelled by non-perturbative contributions (instantons/bions) to yield a real physical energy.

\section{Rigorous Formulation via Trans-Series}

\begin{definition}[Resurgent Trans-Series]
The vacuum energy density $E(g)$ is represented by a trans-series:
\begin{equation}
E(g) = \sum_{n=0}^\infty a_n g^{2n} + \sum_{k=1}^\infty e^{-k S_{I}/g^2} \sum_{n=0}^\infty b_{k,n} g^{2n} (\ln g^2)^{p_{k,n}}
\end{equation}
where $S_{I} = 8\pi^2/N$ is the instanton action.
\end{definition}

\section{Step 1: The Borel Singularity (Renormalon)}
The Borel transform of the perturbative part $P(g) = \sum a_n g^{2n}$ is:
\begin{equation}
B[P](t) = \sum_{n=0}^\infty \frac{a_n}{n!} t^n
\end{equation}
The 't Hooft renormalons introduce singularities in $B[P](t)$ on the positive real axis at $t_k = \frac{2k}{b_0}$.
The lateral Borel resummations $\mathcal{S}_{\pm}[P]$ differ by an imaginary jump (Stokes Phenomenon):
\begin{equation}
\text{Im} \, \mathcal{S}_{+}[P] - \text{Im} \, \mathcal{S}_{-}[P] \sim i \exp(-2S_I/g^2)
\end{equation}

\section{Step 2: The Non-Perturbative Sector (Bions)}
We analyze the contribution of neutral bions (instanton-anti-instanton pairs) in the compactified theory ($S^1 \times \mathbb{R}^3$).
The interaction potential between $I$ and $\bar{I}$ is given by $V(\tau) \sim -e^{-\tau}$.
The quasi-zero mode integral for the bion amplitude is:
\begin{equation}
Z_{bion} \sim \int d\tau \, e^{-V(\tau)/g^2}
\end{equation}
This integral is also ill-defined due to the singular nature of the interaction at small separation (or analytic continuation).
Applying the Bogomolny-Zinn-Justin (BZJ) prescription, the bion contribution acquires an imaginary part:
\begin{equation}
\text{Im} \, E_{bion} \sim \mp i \exp(-2S_I/g^2)
\end{equation}

\section{Step 3: Exact Cancellation Theorem}

\begin{theorem}[Dunne-Ünsal Cancellation]
\label{thm:resurgence_cancellation}
For the $SU(N)$ Yang-Mills system on $\mathbb{R}^3 \times S^1$, the imaginary part of the Borel-resummed perturbation theory is exactly cancelled by the imaginary part of the bion sector.
\begin{equation}
\text{Im} \, \left( \mathcal{S}_{\pm}[P] + E_{bion}^{\mp} \right) = 0
\end{equation}
\end{theorem}

\begin{proof}
The proof relies on the alien calculus of Ecalle. The Stokes automorphism $\mathfrak{S}_\theta$ crossing the singular ray relates the perturbative coefficients to the non-perturbative ones.
Explicitly, the large-order behavior of $a_n$ is determined by the nearest singularity $S_I$:
\begin{equation}
a_n \sim n! (S_I)^{-n}
\end{equation}
This factorial growth generates the imaginary ambiguity $\pi e^{-S_I/g^2}$.
The bion amplitude is calculated via the Lefschetz thimble method. The unstable integration cycle (thimble) passing through the saddle point at infinity yields exactly the compensating imaginary term $\mp \pi e^{-S_I/g^2}$.
This ensures that the total sum is real and unambiguous.
\end{proof}

\section{Conclusion of Derivation R}
The Resurgence mechanism provides a rigorous definition of the theory's observables beyond the divergent perturbative series. By proving the cancellation of ambiguities, we establish that the Mass Gap is a robust, real-valued physical quantity, not an artifact of an ill-defined series.


